
définissant \(Z\), et soit \(z\) le point générique de \(Z\). Comme \(Z\) est intègre, \(\underline{J}_z\) est l'idéal maximal de \(\underline{O}_{X,z}\), et, puisque \(X\) est régulier, \(\underline{J}_z\) peut être engendré par une suite \(\underline{O}_{X,z}\)-régulière. On peut donc trouver une suite de sections de \(\underline{J}\) dans un voisinage de \(z\), formant une suite \(\underline{O}_X\)-régulière. Par suite, il existe un voisinage ouvert \(U\) de \(z\), tel que \(Z'=X-U\) soit contenu dans \(Z\), et que \(\underline{J}|U\) soit un idéal régulier de \(\underline{O}_U\). L'immersion \(Z\cap U\xrightarrow{i}U\) est alors une immersion régulière, dont la codimension est \(\dim(\underline{O}_{X,z})=\codim(Z,X)\) (EGA IV 5.1.2). Par ailleurs, on a la suite exacte (IV\qquad) :
\[
  K^{\bullet}(Z')\xrightarrow{j_*}K^{\bullet}(X)
  \xrightarrow{u^*}K^{\bullet}(U)\to0,
\]
où \(j\) est l'immersion \(Z'\to X\) et \(u\) l'immersion \(U\to X\) ; on obtient une suite analogue après tensorisation par \(\mathbb Q\). Or \(u^*(\underline{O}_Z)=i_*(\underline{O}_{Z\cap U})\), et d'après 4.6, l'image \(u^*(1_Z)\) de la classe de \(\underline{O}_Z\) par \(u^*\) est donc dans \(\Fil^kK^{\bullet}(U)_{\mathbb Q}\). Comme \(u^*\) est un \(\lambda\)-homomorphisme surjectif, il induit une \emph{surjection} \(\Fil^kK^{\bullet}(X)\longrightarrow \Fil^kK^{\bullet}(U)\), donc \(u^*(1_Z)=u^*(x)\) avec \(x\in\Fil^kK^{\bullet}(X)\) ; il existe donc \(y\in j_*(K^{\bullet}(Z'))\) tel que :
\[
  \operatorname{cl}(\underline{O}_Z)-y\in\Fil^kK^{\bullet}(X)_{\mathbb Q}.
\]
Or \(Z'\subset Z\), et comme \(z\notin Z'\), \(\codim(Z',X)>k\). Par suite, \(y\) peut s'exprimer comme combinaison de classes de faisceaux \(\underline{O}_{Z''_\alpha}\), où les \(Z''_\alpha\) sont des sous-schémas fermés intègres de \(Z\), distincts de \(Z\). D'après l'hypothèse de récurrence, \(y\in\Fil^{k+1}K^{\bullet}(X)_{\mathbb Q}\), ce qui achève la démonstration.

\subsection*{BIBLIOGRAPHIE}
\begin{enumerate}[label={[\arabic*]}]
\item Cartan--Eilenberg, Homological Algebra.
\end{enumerate}



% Expose VIII macro additions
\providecommand{\etale}{\'etale}
\providecommand{\Ch}{\operatorname{Ch}}


